\begin{abstract}
O.~Caramello's \q{bridges} read one property of a topos on two presentations of it and turn the two readings into a theorem relating the presentations. A companion article \cite{bridges-generator} runs such bridges by hand as a theorem generator. This note applies the same generator to Grothendieck's manuscript notes (Archives Grothendieck, Université de Montpellier). Three arches (two-valued, connected, presheaf topos of a preorder), proved in Lean~4 in this edition's package \texttt{lean/bridges}, are read on three decks that recur in his notes: the simplex category, groups and preorders. We discuss two precedents in his notes: an exercise of folder 161-3, p.~26, which the generator re-derives from the Morita arch and whose preorder form is proved in Lean, with the same identification in folder 37, p.~48; and the correspondence between subtopoi and quotients of a theory in folder 161-2, pp.~17 and~44. A last deck from folder 161-2, the topos of pseudo-fields (von Neumann regular rings), gets a section of its own: we check one of his three presentations of it, read the arches on it and answer a question he left in the margin. An appendix lists the Lean files. A French version accompanies this text.
\end{abstract}
